\chapter{Hypoxemia and Hypoxia}

\section*{Definition}
Hypoxemia is abnormally low oxygen in arterial blood; hypoxia means inadequate oxygen delivery or use by tissues. Severe oxygen deficiency can rapidly damage the brain, heart, and other organs.

\section*{What Happens in Your Body}
The presentation reflects failure of normal airway, respiratory, circulatory, neurologic, immune, gastrointestinal, muscular, or renal function. Rapid assessment focuses on immediate threats to airway, breathing, circulation, neurologic status, and reversible metabolic causes.

\section*{Causes}
\begin{itemize}
  \item Pneumonia, asthma, COPD, pulmonary embolism, pneumothorax, pulmonary edema, heart failure, airway obstruction, anemia, shock, high altitude, and impaired ventilation.
  \item Medication effects, procedures, toxins, trauma, and underlying chronic disease may contribute.
  \item Severe presentations can deteriorate quickly even when early physical findings are limited.
\end{itemize}

\section*{When to See a Doctor}
Severe breathlessness, blue or gray lips, confusion, inability to speak normally, chest pain, fainting, extreme drowsiness, or a low oxygen reading with symptoms requires emergency care. Do not delay emergency assessment while attempting home treatment.

\section*{Related Symptoms}
\begin{itemize}
  \item Fever, pain, weakness, dizziness, nausea, vomiting, or change in normal function
  \item Breathlessness, cyanosis, low blood pressure, reduced urine, bleeding, or confusion in severe illness
\end{itemize}

\section*{Self-Care / Home Management}
These presentations should not be managed independently when red flags are present. Follow an established emergency action plan if one exists, bring medication and medical history information, and do not give food, drink, or oral medicine to a person who is unconscious or unable to swallow safely.

\section*{How Your Doctor Will Evaluate You}
\begin{itemize}
  \item Immediate assessment includes airway, breathing, circulation, mental status, vital signs, oxygen saturation, blood glucose, and focused examination.
  \item Blood count, electrolytes, renal and liver tests, blood gas, cultures, ECG, imaging, and targeted toxicology or immune testing may be required.
  \item The history covers onset, exposures, medications, cancer treatment, infection risk, trauma, comorbidities, and baseline function.
\end{itemize}

\section*{Treatment Options}
Treatment is cause-specific and may include emergency airway and breathing support, intramuscular epinephrine for anaphylaxis, intravenous fluids, antibiotics, glucose, antidotes, seizure treatment, blood products, renal replacement therapy, or specialist intervention. Monitoring and reassessment are essential.

\section*{References}
\begin{thebibliography}{99}
\bibitem{symptom_hypoxemia:ref1} Jameson JL, et al. \textit{Harrison's Principles of Internal Medicine}. 21st ed. McGraw-Hill; 2022.
\bibitem{symptom_hypoxemia:ref2} American Heart Association and relevant specialty societies. Current emergency evaluation and management guidance; accessed 2025.
\bibitem{symptom_hypoxemia:ref3} National Institute for Health and Care Excellence. Relevant clinical guidance. NICE; accessed 2025.
\end{thebibliography}
